\chapter{Deep Vein Thrombosis Symptoms}

\section*{Definition}
Deep vein thrombosis (DVT) is the formation of a blood clot within a deep vein, typically in the lower extremities. Clinical diagnosis uses the Wells score to stratify pre-test probability; D-dimer testing and compression ultrasound confirm the diagnosis.

\section*{What Happens in Your Body}
Virchow's triad (venous stasis, endothelial injury, hypercoagulability) promotes thrombus formation within deep veins. The clot occludes venous return, causing distal edema, pain, and inflammation; embolization of the thrombus to the pulmonary circulation may cause pulmonary embolism.

\section*{Causes}
\begin{itemize}
  \item \textbf{Stasis:} Prolonged immobilization (long-haul travel, hospitalization, surgery), obesity, paralysis, heart failure
  \item \textbf{Endothelial injury:} Trauma, surgery (especially hip/knee arthroplasty), central venous catheter, vasculitis
  \item \textbf{Hypercoagulable states:} Factor V Leiden mutation, prothrombin G20210A mutation, antiphospholipid syndrome, protein C/S deficiency, antithrombin deficiency, malignancy, pregnancy, oral contraceptives/hormone therapy, polycythemia vera
  \item \textbf{Other:} Age greater than 60, prior DVT/PE, thrombocytosis, nephrotic syndrome, inflammatory bowel disease
\end{itemize}

\section*{When to See a Doctor}
See your doctor if:
\begin{itemize}
  \item Unilateral calf or thigh swelling, pain, warmth, or erythema
  \item Sudden onset of dyspnea, chest pain, hemoptysis, or syncope (suggestive of concurrent PE)
  \item Recent surgery, immobility, or known hypercoagulable state with leg symptoms
  \item Palpable cord-like vein, edema (increased limb circumference), or pitting edema
  \item Homan's sign (calf pain on dorsiflexion) is neither sensitive nor specific enough for diagnosis
\end{itemize}

\section*{Related Symptoms}
\begin{itemize}
  \item Unilateral leg swelling, pain, warmth, erythema, dilated superficial veins, calf tenderness, low-grade fever
\end{itemize}
\textbf{Important:} Up to 30\% of DVTs are asymptomatic, and pulmonary embolism may be the first clinical manifestation; the absence of leg symptoms does not exclude DVT.

\section*{Self-Care / Home Management}
\begin{itemize}
  \item Ambulation as tolerated once anticoagulation is initiated; do not immobilize the leg
  \item Elevate the affected leg above heart level when resting to reduce edema
  \item Apply warm compresses to relieve local pain and inflammation
  \item Wear compression stockings (30--40 mmHg) to reduce post-thrombotic syndrome
  \item Strictly adhere to anticoagulation regimen; do not skip doses
\end{itemize}

\section*{How Your Doctor Will Evaluate You}
\begin{itemize}
  \item Wells score calculation: active cancer (+1), paralysis/paresis (+1), recent immobilization/surgery (+1), localized tenderness (+1), entire leg swollen (+1), calf swelling greater than 3 cm (+1), pitting edema (+1), alternative diagnosis as likely or more likely (-2)
  \item D-dimer (ELISA): highly sensitive (greater than 95\%) but non-specific; negative test rules out DVT in low-to-moderate probability
  \item Compression venous duplex ultrasound: non-compressibility of the vein is diagnostic; sensitivity and specificity greater than 95\% for proximal DVT
  \item Venography or CT venography reserved for equivocal cases or when ultrasound is technically limited
\end{itemize}

\section*{Treatment Options}
\begin{itemize}
  \item Anticoagulation: low-molecular-weight heparin (LMWH) or fondaparinux first-line, transitioning to direct oral anticoagulants (rivaroxaban, apixaban) or warfarin
  \item Duration: minimum 3 months for provoked DVT; indefinite for unprovoked or recurrent DVT
  \item Thrombolysis reserved for massive DVT with phlegmasia cerulea dolens or limb-threatening ischemia
  \item Inferior vena cava filter for people with contraindications to anticoagulation or recurrent PE despite anticoagulation
  \item Compression stockings (30--40 mmHg) for 2 years to prevent post-thrombotic syndrome
\end{itemize}

\section*{References}
\begin{thebibliography}{99}
\bibitem{pe1} Konstantinides SV, Meyer G, Becattini C, et al. "2019 ESC guidelines for the diagnosis and management of acute pulmonary embolism." Eur Heart J. 2020;41(4):543-603.
\bibitem{pe2} Di Nisio M, van Es N, Buller HR. "Deep vein thrombosis and pulmonary embolism." Lancet. 2016;388(10063):3060-3073.
\bibitem{pe3} Thompson BT, Kabrhel C. "Overview of acute pulmonary embolism in adults." UpToDate, 2023.
\bibitem{pe4} Kearon C, Akl EA, Ornelas J, et al. "Antithrombotic therapy for VTE disease: CHEST guideline." Chest. 2016;149(2):315-352.
\bibitem{pe5} Wells PS, Anderson DR, Rodger M, et al. "Derivation of a simple clinical model to categorize patients probability of pulmonary embolism." Thromb Haemost. 2000;83(3):416-420.
\bibitem{pe6} Becattini C, Agnelli G. "Acute pulmonary embolism." In: Harrison\'s Principles of Internal Medicine. 21st ed. McGraw-Hill; 2022.
\end{thebibliography}
